% chapters/ch05.tex —— 第 5 章
% 体例：具体失败 → 设计问题 → 机制解释 → 对照实验 → 条件变化 → 选择原则
% 数字身份纪律：A-Mem / Trap of Trajectory / KSI 的数字一律 \cited{}（出处见 source-map 第 5 章各条）；
% E03 / E03c 的本书实测一律 \measured{}；由实测算出的量用 \calcnum{}。
\chapter{记忆如何转化为可复用的经验}
\label{ch05}

\chaptergoal{保存更多历史，为什么不一定更会做事？}
\chapterprops{P3（形式与容量如何共同影响迁移）与
P3b（拟合目标怎样影响迁移）。本章用E03与E03c的观察解释表示、读取者与比较范围之间的关系。}
\chapterfig{图~\ref{fig:ch05-fit}：一条经验可以拟合到求解器，也可以拟合到问题域}
\chapterpitfall{把``存下来了''当成``学会了''；把检索到的相关性当成因果；
把同族新题上的提升当成跨族迁移}

% ---------- 1. 具体失败 ----------
\begin{failurebox}
第 4 章的栈错误，在另一项实验中又成为一条经验来源。
E03 拆分 A 的同一份失败记录，被渲染成三种形式：截取记录、短摘要、结构化经验。
短摘要明确指出需要删除并返回栈顶；结构化条目进一步说明，查看栈顶的场景不应照搬这个操作。
名为 \texttt{form\_raw} 的条目却在 \measured{600} 个字符处硬截断，
代码停在 \texttt{push} 方法的一部分，还没有出现出错的 \texttt{pop} 方法。

“保留原始记录”这个名字，并不表示关键信息完整保留下来。
与此同时，不同渲染改变了检索排名：拆分 A 的 \measured{16} 道留出题中，
结构化与摘要选择的来源条目集合有 \calcnum{12} 道不同；截取记录与摘要有 \calcnum{14} 道不同。
这些差异已与实际系统提示逐项核对。

所以本章不能只问哪一种记忆文字更好。还要问：
什么被写下、什么被截掉、什么被取出，以及读到它的系统最终做了什么。
\end{failurebox}

\designquestion{判据不是``记住了多少''，而是：
\textbf{这条经验在换了一件事之后还成不成立}。
换一道题、换一个族、换一个读到它的模型——
每一种变化都可能改变经验的适用性，这一步必须可检查。}

\section{同一段经历，为什么会写成三种不同的经验}
\label{sec:ch05-same-source}

同源材料是一个重要控制，却只是控制的起点。E03 先让无记忆 Agent 产生来源记录，再从同一条记录渲染不同形式。这样避免了为摘要选一批容易题、为结构化经验另选一批难题。渲染完成以后，内容覆盖、长度、措辞和检索排名仍可能分化。若把“同源”直接解释成“其他都相同”，后面的归因就会跳过这些实际变化。

\subsection{栈记录中保留与丢失的内容}

拆分 A 的栈来源记录包含失败验收和最后代码。归档的短摘要说，应同时返回并移除栈顶，可以使用容器的弹出方法，不能只读末元素。结构化版本则分成事实、适用、不适用三行：事实指出只返回未删除；适用范围是实现需要改变内容的弹出操作；不适用范围是只查看栈顶而不能修改内容的操作。

这两份文本都比原始记录短，却保留了诊断的关键。它们之间的区别也不只是字段名：结构化版本明确写了相反需求，短摘要更直接给出代码操作建议。若后续行为不同，可能来自条件边界，也可能来自具体操作提示、措辞长度或读取者对格式的反应。要分离这些因素，需要设计新的受控版本，旧运行没有同时完成所有消融。

截取记录臂提供了另一种取舍。它按固定顺序拼接源题、工具名称、验收状态、部分验收输出与代码，然后取前 \measured{600} 个字符。在这条失败记录里，错误回溯和临时文件路径占据了前部空间，末尾代码被截到 \texttt{push} 的一部分。恰恰需要检查的 \texttt{pop} 函数没有进入条目。

因此，“原始轨迹”在这里是历史臂名，不是逐步完整档案。条目中没有完整保留每次工具参数、每段模型回答或全部观察，连拼接后的记录也会截断。完整任务轨迹仍可用于审计，但被注入下游上下文的是这个短片段。评价经验形式时，应描述实际展示给读取者的材料，不能用归档材料的完整性替代提示内容的完整性。

这个例子也说明，固定字符上限并不保证各条经验受到相同程度的压缩。成功输出通常短，失败回溯可能长；同样的拼接顺序会使失败记录留下更少代码。如果长期按此规则处理，保留内容可能系统性地随结果类别变化。这里提出的是可检查的设计风险，不是未经统计就断言所有失败记录都会丢失关键内容。

\subsection{长度口径决定了“同容量”究竟是什么意思}

拆分 A 每种形式各有 \measured{8} 条条目，平均长度分别为截取记录 \measured{600}、摘要 \measured{44}、结构化 \measured{102.75} 个字符。这里的字符按 Python 字符串长度计数，包含中文、英文、数字、空白与标点，不是汉字数，也不是计费 token 数。拆分 B 的求解器文本与问题域文本平均长度分别为 \measured{60.75} 和 \measured{104.375} 个字符。

相同条目数控制了存储单元数量，没有控制上下文长度。即使每次都取两条，注入的字数和 token 数仍不同；即使把字符数截成相同，模型分词后的 token 数也可能不同。容量可以指条目数、磁盘字节、索引大小、检索条数或提示预算，必须选定所研究的那一种。

长度还不是内容量的充分代理。重复模板、路径和测试框架文字占字符，却未必提供当前任务需要的信息；一个短条件可能排除大量错误操作。反过来，删掉具体输入只留下抽象口号，也可能使诊断失去依据。因此，应把长度当作资源约束，把信息覆盖作为另一个需要检查的对象，而不是认为压得更短就必然更好。

原提示要求摘要不超过一段较短文字，但实现还有自己的字符截断。提示中的限长是一项请求，实际输出长度才是观测。归档中确有摘要条目超过提示宣称的字符数量，不能把提示目标直接写成严格执行的上限。对可复现实验，存储前后的长度与最终展示长度应分别可查。

\subsection{改写经验也改变了哪些经验被看到}

E03 沿用字符二元组检索。去掉空白后，将查询和条目分别变成相邻字符对的集合，按查询与条目共同字符对的数量除以条目字符对数量排序。这个分母来自候选条目，不是两个集合的并集，所以它不是常见的 Jaccard 系数。长文本若增加了许多不与题目重合的字符对，分数可能下降；较短文本若保留了几个匹配词，分数可能上升。同分时，保存列表中位置较后的条目优先。因此，复算检索结果还需要保留条目顺序。

这不是说长文本总被惩罚，或者短文本总被优待。新增词语也可能提高交集，具体排序由两部分共同决定。但格式字段、代码标识、自然语言同义表达都会改变集合，故“只改存储形式、检索器没改”仍不等于检索结果相同。检索器的函数固定，输入表示变化也会改变输出。

离线复核使用归档经验与实际任务陈述重建排名，并将选中文字与原始系统提示逐项比对。拆分 A 的三种记忆臂共核对 \calcnum{240} 个提示，拆分 B 的两种同模型记忆臂共核对 \calcnum{160} 个提示，选中文字均出现于对应提示。拆分 A 中，结构化与摘要的源条目集合在 \calcnum{12/16} 道题上不同，截取记录与摘要在 \calcnum{14/16} 道题上不同；拆分 B 的两种拟合目标也有 \calcnum{12/16} 道题不同。

例如标准拆分中的 \texttt{classstate-07}，求解器文本选中了异常契约题与栈题的经验，问题域文本选中了计数器题与可变默认参数题的经验。仅这一处就足以说明，读取者看到的不是相同两条事实的不同排版。它还不能证明哪一组更适用，更不能证明哪条经验导致成败；那需要继续看任务契约、实际行动和反事实对照。

\subsection{把检索变化视为机制，还是视为混淆}

答案取决于问题。如果研究的是“把整个经验系统从摘要切换到结构化格式会怎样”，检索变化就是完整方案效果的一部分，应保留并报告。如果研究的是“在相同经验内容已经被选中时，加上条件字段是否改善执行”，那么检索变化就是需要控制的中间路径。

后一种设计可以先固定来源条目标识，再渲染各版本给读取者。这样切断了形式经由检索排名影响成绩的路径，却也偏离了实际自动检索流程。前一种设计更接近部署，后一种更适合拆解机制。两种结果可以并列，但不能用一种设计的区间回答另一种设计的问题。

还可以增加一个中间检查：固定检索器与经验来源，记录每道题选中的来源集合、顺序和展示长度。若整体成绩变化，这些记录有助于定位下一步该做哪项干预。它们仍属于机制观察，不能仅凭“更相关的经验出现得更多”就宣称中介效应已经被识别。

\subsection{先确认信息还在，再检查能否取出}
\label{sec:ch05-construction-retrieval}

同一个错误答案可以来自不同环节。以章首栈记录为例：若条目已丢掉“移除栈顶”的要求，单换检索器不能恢复这个条件；若条件仍在库里，却没有进入提示，应先查候选选择；若正确条件已经展示，执行者仍只读取末元素，就要查它怎样理解并使用经验。这些路径可以并存，需要逐题沿保存内容、检索结果和实际动作核对。

Zhou 等在 LongMemEval-S 与 BEAM-100K 上比较六种显式记忆和原始历史 RAG，固定已构建的记忆，再改变检索器与上下文预算。显式构建没有取得稳定的总体优势，收益集中于部分题型和受限条件\cite{raw-history-memory-2610}。其控制保留了各系统周边的查询规划、链接扩展等流程，不能当作只改变文本表示的纯形式实验。

该研究还按标注的支持语句追溯来源，构造不受实际检索排名约束的参照\cite{raw-history-memory-2610}。这有助于定位检索瓶颈，却不是可部署的读取策略；来源范围重合，也不保证改写后的条目完整保留答案。若直接提供候选支持材料后仍答错，还应检查表达与读取能力，不能仅凭答案错误就断言信息已经被删除。

\section{从一次经历抽象出规则，哪些细节可以删除}
\label{sec:ch05-abstraction}

抽象的目标不是让经验听起来更普遍，而是在保留必要条件的前提下去掉偶然细节。需要删除的可能是临时路径、无关对象名与重复日志；需要保留的可能是异常契约、状态是否共享、输入能否修改，以及失败出现的具体边界。对每项细节都可以问：如果它改变，建议是否还成立？

\subsection{按变量作用区分必要条件与背景}

以下是本书构造的教学例子。某次文件处理任务要求读入数据、按指定键排序并输出新文件，原文件必须保留。失败发生在输出覆盖原文件之后。路径里的日期、用户名和临时目录通常属于背景；“输出与输入不能是同一目标”则直接连接到失败条件。若把经验写成“不要用昨天的路径”，保存的是一次偶然关联；写成“写入前核对目标与受保护输入是否相同”，才指向可检验的机制。

抽象仍不能过宽。某些任务本来就要求原地转换，并有备份和恢复保证；一律禁止同路径写入会改变需求。更合适的条件是“在输入必须保持不变的任务中”，再给出检查动作。反例不是为了把经验写得冗长，而是防止一条正确的局部规则抢占所有场景的决策权。

判断某个细节是否必要，可以构造最小变化：只改目录名、只改输入内容、只改是否允许覆盖，再检查原建议是否仍然合理。这些变化可以先由规格分析完成；需要程序行为支持时，再做受控执行。若同时改变任务目标与运行环境，观察到不同结果也无法知道哪项条件真正重要。

\subsection{从成功中提炼经验，也要保留不确定性}

成功轨迹往往包含很多不必要动作。系统先读了一份无关文档，后来凭已有知识完成任务；若摘要把整个顺序都写成成功步骤，就会把无关读取固化为前置条件。下一题照做这些步骤，可能只是增加成本，也可能引入错误信息。

因此，成功经验可以分成已观察事实与候选策略。事实是这次使用某个版本、完成某个结果；策略是以后遇到某条件可以尝试某动作。后者需要更多证据，尤其要检查省掉这个动作是否仍然成功。没有反事实比较时，可以说“本次有效流程的一部分”，而不应说“成功所必需”。

失败经验同样有这个问题。一次失败后换了多个设置才成功，不能把最后一次改动自动认定为原因。把中间候选及检查结果保留下来，可以更准确地界定目前知道什么；无法分离的联合修改，就先作为联合策略记录。过早概括为单一规律，会让经验更容易检索，却更难正确使用。

\subsection{事实、解释与建议不要互相冒充}

一个条目可以写“该输入下返回空集合”，也可以写“因为默认值共享导致状态丢失”，还可以写“改用独立容器”。第一句是观察，第二句是解释，第三句是干预建议。它们需要不同的支持材料。观察可以由日志核对，解释需要与其他可能原因竞争，建议需要检查修改后的后果。

若把三句压缩成“默认容器错误”，下游既不知道出错的是共享还是清空，也不知道返回空集合究竟违反了什么要求。简短到无法恢复推理对象的经验，可能只剩一个关键词。结构化并非必须使用固定字段名，但应让这些逻辑角色可区分，否则一句推测很容易在反复转述中变成“已证事实”。

KSI 的条目模式把适用条件、不适用条件与证据并列\cite{ksi-2607}，提供了一种实现方式。本书借用的是可检查边界与来源这一思想，不把特定字段名当通用正确性保证。若字段里只是“适用于类似情况”“不适用于不同情况”，形式合格，信息仍然不足。

\subsection{保留一个能推翻规则的测试}

经验条目不必随身携带庞大测试集，但最好指出什么观测会使它失效。例如“输入必须不变”可以由调用前后比较检查；“这个接口返回新对象”可以由对象身份与内容分别检查；“该操作能安全重试”则需要检验重复调用是否产生额外副作用。

这种反驳条件能把经验从口号变成可维护的资产。新版本出现反例时，系统可以判断是条件范围过宽、环境假设过期，还是原规则本来就不成立。若只保存“通常这样更好”，几乎任何失败都能被解释成例外，经验便很难真正被修订。

并非所有知识都能立即形式化。对于写作风格或复杂业务判断，可以保留任务样例、用户选择和争议说明。关键是不要伪装成更强的证据：偏好例子支持特定情境下的选择，不自动成为跨用户的质量定律。抽象要提高可复用性，也要保留我们尚不知道的边界。


\section{四种形态：轨迹、事实、策略与程序}
\label{sec:ch05-forms}

\begin{table}[htbp]
\centering
\caption{经验的四种形态：常见检索线索与典型失效方式}
\label{tab:ch05-forms}
\begin{tabular}{>{\raggedright\arraybackslash}p{2.2cm}>{\raggedright\arraybackslash}p{4.0cm}>{\raggedright\arraybackslash}p{3.6cm}>{\raggedright\arraybackslash}p{3.3cm}}
\toprule
形态 & 一条是什么 & 常见检索线索 & 典型失效 \\
\midrule
轨迹 & 一次完整尝试的原始记录 & 与当前任务的表面相似 & 长、噪声多、含大量与成败无关的步骤 \\
事实 & 关于世界或任务域的一条陈述 & 实体与条件 & 过期；被更晚的事实覆盖 \\
策略 & ``遇到 X 就做 Y''的 conditional & 条件匹配 & 条件写得太宽，于是被套到不该套的地方 \\
程序 & 可执行的一段代码 & 签名与前置条件 & 接口漂移；依赖消失 \\
\bottomrule
\end{tabular}
\end{table}

表~\ref{tab:ch05-forms} 的四行不是并列的优劣排序，而是四种不同的失效方式。
本书第 7 章处理的是``事实与策略''被写进上下文的形态，第 8 章处理``程序''；
本章关心的是它们共同的前置问题：\emph{一条经验凭什么被留下}。

\section{筛选、抽象、合并与冲突处理}
\label{sec:ch05-curate}

A-Mem 把一条记忆写成一个带七个字段的笔记
$m_i=\{c_i,t_i,K_i,G_i,X_i,e_i,L_i\}$：内容、时间戳、关键词、标签、上下文描述、
embedding 与链接集合\cite{amem-2502}。
新笔记写入时按余弦相似度
（假定两个嵌入向量均非零）

\begin{equation}
s_{n,j}=\frac{e_n\cdot e_j}{\lvert e_n\rvert\,\lvert e_j\rvert}
\label{eq:ch05-link}
\end{equation}

取 top-$k$ 近邻，再由模型决定是否建立链接，并对近邻做一次演化\cite{amem-2502}。

\begin{sloppypar}
成本需要连同计量对象一起读。A-Mem 的成本段报告每次记忆操作约 \cited{1{,}200} 个 token，
并汇总相对于约 \cited{16{,}900} token 基线的 \cited{85--93\%} 降幅；
单次低于 \cited{0.0003} 美元是该论文当时的接口成本报告，不是当前报价\cite{amem-2502}。
这两个近似 token 数直接计算只能得到约 \calcnum{92.9\%}，不能独自推出整个降幅区间。

原文表 1 另列回答一道问题的平均 token 长度：GPT-4o-mini 配置为 \cited{2{,}520}，
GPT-4o 配置为 \cited{1{,}216}。这些表格数与成本段的操作口径应分别保留，
不能把约一千二百写成所有模型、所有写入与读取流程的统一费用。
若要比较完整生命周期，还需列出初始化、链接生成、近邻改写、检索和回答各自发生多少次，
以及哪部分已经计入原报告。论文的局部成本优势值得研究，但不自动证明完整方案更省或经验更正确。
\end{sloppypar}

\begin{principlebox}
\textbf{按实际消融读模块贡献。}
A-Mem 的表 3 比较同时去掉链接生成与记忆演化、仅去掉记忆演化、以及完整方案\cite{amem-2502}。
这支持所测试配置下的模块比较，不能说原文完全没有“有内容、无链接”的参照。
但它也不是在其他输入与读取过程逐字不变时，仅开关图边的纯拓扑实验。
若要把收益进一步归因于某种链接结构，需要针对这个问题另设控制。
\end{principlebox}

删除、改写与降权解决不同问题。A-Mem 的方法描述重点是笔记、链接与演化，
不能仅凭没有独立删除模块就推断所有矛盾必然保留。
其他方法即使没有单独的“删错惩罚”，删除也可能通过下游成绩受到间接反馈。
真正需要审查的是：错误条目怎样停止影响行为，误删以后怎样恢复，以及这些结果有没有被测量。

\section{合并经验时，怎样保留证据而不制造共识}
\label{sec:ch05-merge}

同一个问题被多次遇到以后，经验库里会出现许多相似条目。简单追加会增加检索负担，直接合成一句“共同规律”又可能抹掉差异。合并应先辨认重复的是文字、事件还是主张：三条文字相近的经验可能来自三次独立观察，也可能都转述同一次失败；三条措辞不同的经验也可能指向同一个条件。

\subsection{先去重事件，再合并主张}

如果一份失败报告被摘要器、审校器和汇总器分别改写，最后得到三个条目，这仍只有一个来源事件。把条目数当证据数，会使多次转述获得不应有的权重。每条派生经验应指向最初的运行或外部来源，合并时沿来源关系去重，而不只比较最终文本。

反过来，两个独立任务都出现同一种错误，即使最终经验句完全一样，也不应丢掉第二次观察。可以保留一份规范表述，并把不同证据挂在其后，分别记录各自条件。这样减少提示中的重复，又不牺牲证据的覆盖。若两次任务共享同一测试模板或同一环境缺陷，它们的独立性还要进一步限定。

一个实用的记录方式是分开“条目标识”和“主张标识”。条目标识对应某次具体写入，主张标识对应被维护的命题。若某条写入后来被发现是误转述，可以撤销这份支持而不必抹掉其他来源；若主张本身错误，则应沿主张关系检查全部派生建议。这个结构是本书的工程建议，不是现有 E03 归档已经具备的能力。

\subsection{条件不相同，先分支，不急着选边}

考虑两条教学示例：“转换失败时返回默认值”和“转换失败时应抛出错误”。它们看起来冲突，实际上可能来自不同接口契约。一个交互式表单希望继续处理其他字段，另一个底层校验库要求调用方明确处理错误。先补齐调用场景与接口要求，矛盾可能变成两个条件分支。

合并后可以保存共同部分：转换存在失败分支，失败行为必须符合调用契约；两个具体分支分别保留。若直接多数投票，常见场景可能淹没较少见但同样正确的契约。经验数量多，不代表它有权覆盖当前任务的明确要求。

也有真正冲突：两个条目针对同一版本、同一输入范围、同一目标，却分别声称某操作会与不会修改输入。此时需要回到来源或实际检查，而不是加一句“视情况而定”。条件已经相同，继续模糊条件只会回避问题。若证据暂时不足，保留争议状态并限制自动使用，比任意择一更可追溯。

\subsection{合并不能把窄证据扩成宽命题}

两条经验分别说某函数在空列表和单元素列表上正确，合并以后不能写成“适用于任意长度列表”。被观察条件的并集只有这两个范围，不是所有可能输入。类似地，两个模型都能读取某条经验，不代表它对任何模型可移植；两个时间点都有效，也不证明期间每个版本都有效。

经验摘要常会为了简洁删掉限定词，合并比单次摘要更容易发生这种扩大。可以在合并前后比较主张范围：新版允许哪些原来未允许的输入、动作或环境？如果出现新增范围，必须把它标成待验证扩展，而不是默认为已有证据的一部分。

合并也可以缩小范围。若旧条目声称所有可变对象都应该复制，后续发现某个接口明确要求原地修改，就应收窄到需要保持调用方对象不变的情况。这不是“记忆退化”，而是知识边界变得更准确。效果评估仍需检查收窄以后是否影响目标任务，但不能只按条目变短或变长判断好坏。

\subsection{版本关系比覆盖写入更适合解释变化}

用新文字覆盖旧文字很方便，却会使历史行动无法解释。某次任务遵循的是旧条件，审计时只看到已经修正的新条件，便可能误判执行者违背了当时规则。更好的做法是保留版本关系：新版取代哪一版，改变了哪些条件，依据哪份证据，何时对读取者生效。

读取时通常只展示当前有效版本，审计时仍能回到历史版本。这样不必把旧错误持续放进提示，也不必删除形成错误的过程。版本保留与提示可见性是两项不同决定，不能用“需要审计”作为永远启用旧经验的理由。

CALM 提供了一个将两项决定分开的实例。在其预定义的关系规则下，强冲突更新把新记录置为 \texttt{ACTIVE}、旧记录置为 \texttt{SUPERSEDED}，同时保留替换指针与有效时间边界\cite{calm-2610}。明确询问过去的查询，可让后端提供且满足内容、时间条件的查询局部副本进入本次使用筛选；旧记录的持久状态保持不变\cite{calm-2610}。这些规则规定记录何时可被采用；状态转移本身不能替代对争议事实的核查。

若条目被许多其他经验引用，合并还涉及依赖迁移。新主张是否保留了被引用的条件？原引用是支持例子、反例还是实现来源？统一重定向链接可能把反对意见变成支持意见。维护关系时应保留边的含义，而不只是让链接不报错。

\subsection{链接的价值在于可追踪关系，不是边越多越好}

A-Mem 通过近邻检索和模型分析建立关联，再更新相关笔记的描述\cite{amem-2502}。这使新经验能够改变旧经验的组织方式，但也把链接质量交给了具体表示与模型判断。语义接近的两条记录可能是支持、反例、同一来源的重复，或只共享任务名；把这些关系都视作相互强化，会丢失重要区别。

在工程实现里，可以为链接保留简单类型与来源。支持关系需要说明支持的是哪部分主张，冲突关系需要说明条件重合在哪里，派生关系需要说明转换步骤。若类型由模型生成，仍需接受检查和修订。类型越细，维护成本越高，不能把复杂图结构本身当作效果证据。

评价链接作用时，可比较相同内容下不同读取策略，也可比较完整链接维护方案与简单检索。前者帮助定位结构贡献，后者衡量实际系统价值。无论选择哪种，都应记录新增构建与维护成本。一个没有带来可测收益的复杂关系网，仍可能适合人工审计，但不应被直接称为更强的自主学习能力。

\section{经验生命周期：候选、启用、限制与撤销}
\label{sec:ch05-lifecycle}

经验保存以后，并不只有“存在”与“不存在”两种状态。刚从一次失败生成的建议，证据强度与经过多种任务检查的规则不同；某个版本失效的程序，也可能在旧环境里仍然有效。用生命周期状态表达这些差异，可以让读取者知道一条经验获得了多大的使用权限。

以下状态设计用于说明维护逻辑，不是要求所有系统照搬固定枚举。最小系统可以更简单，但要能区分尚待核查的候选、当前允许使用的条目、已知范围收窄的条目，以及不再进入自动检索的撤销条目。状态转移应由明确证据或管理决定触发，不能仅靠文字中的高置信语气。

\subsection{候选经验不应自动继承执行权限}

一次成功可能产生有价值的线索，也可能产生对无关步骤的赞美。候选区的作用是让线索可保存、可检查，同时限制它立即支配后续行动。对于低风险的检查提醒，可以允许以建议形式呈现；对于会执行写入、删除或网络动作的程序，准入需要覆盖前置条件、权限与回退。

把候选与启用分开，并不必然意味着每条经验都要人工审批。规则检查、离线回放、独立测试和预先授权的策略都可以承担不同环节。关键是门槛与后果相称，并且来源能够追溯。一个提醒“检查空输入”的条目，与一个自动修复所有空输入分支的脚本，不应因为共享同一句经验而获得相同权限。

准入失败也有保存价值。例如某条建议在本题修复成功，却损害另一项接口要求，可以作为反例保留。反例应进入合适的读取位置，避免被一般相似度检索误当正向规则。记录“尝试过且不成立”，有时比再生成一条近似建议更节省资源。

\subsection{启用状态应绑定环境与读取方式}

同一条自然语言经验，作为提示中的弱提醒与作为工具执行前的强制规则，作用不同。条目通过哪种读取方式的验证，就应限定在哪种方式下使用。把验证过的提醒直接升级成自动阻断条件，会改变系统行为，而不只是改变界面显示。

环境绑定也不能只写一个软件版本号。经验可能依赖返回格式、权限、数据类型或任务目标，其中一些变化不会触发版本号更新。记录可以优先保留与主张直接相关的条件：接口返回空值还是错误、容器是否可变、调用能否重试。只有需要时才扩展到完整环境快照，避免元数据越来越多却没有帮助判断。

读取者变化以后，启用状态需要重新检查。模型能解析三个字段，不代表它会遵守不适用条件；工具能加载一段程序，也不代表当前依赖满足。兼容性检查回答“能不能读”，行为检查回答“是否按预期用”，迁移评估回答“总体是否有益”，三者不应合并成一次加载成功。

\subsection{限制使用往往比立刻删除更准确}

新反例可能只影响一部分条件。例如经验要求在缺省参数下新建列表，但系统后来加入需要跨调用共享的缓存接口。此时可以排除缓存接口，而不是宣布经验全部错误。范围收窄有助于保留原先有效的部分，也能让后续读者知道为何存在例外。

若错误来自过期事实，则可以把旧条目限定到历史时间范围。例如某个接口曾经支持一项参数，后来移除；历史任务审计仍可能需要旧信息，当前调用则不应使用。时间适用范围与存储保留期不同，不能因为旧事实仍有历史价值就继续让它参与当前决策。

当无法可靠表达边界，或者误用代价太高时，可以先停止自动检索，将条目保留供检查。停止使用不表示我们已经找到全部原因；它是一项受证据不足约束的管理动作。应记录触发事件与恢复要求，避免条目在下一次重建索引时又无意恢复。

\subsection{撤销要覆盖派生物与缓存}

删掉经验正文，只能保证这个文件消失。它可能已经被写进提示模板、摘要、嵌入索引、技能说明或其他条目的合并版本。若这些派生物继续存在，系统仍可能使用同一错误建议。撤销需要沿来源关系定位影响范围，并检查哪些副本应删除、重建或标为失效。

并不是所有派生物都必须一起废弃。一个摘要可能同时引用多份独立证据，撤销其中一份后，其他证据仍然支持较窄主张。此时应重新评估支持关系，而不是机械级联删除。缺乏来源链时，这种精细处理很难完成，只能采用更保守的停用或重建策略。

缓存也需要版本意识。索引重建成功并不保证正在运行的执行上下文已经刷新；执行上下文可能继续携带旧提示。可以在任务边界核对经验版本，必要时中止会受影响的操作或重新构造上下文。处理方式取决于风险与系统能力，但不能把存储层的删除时间直接当作所有执行者停止使用的时间。

LEDGER v1 在其合成编码负载中，分别检查送给回答模型的证据和最终答案是否仍含被删除的值\cite{ledger-2611}。这两个检查回答不同问题：答案里没有出现该值，并不能证明它没有进入模型输入。论文列明，删除操作不覆盖旧快照、响应缓存和运行日志；这些副本需要另行规定保留与清理规则。

\subsection{遗忘策略需要看机会成本}

经验库容量有限时，维护者必须决定留下什么。按最近使用淘汰容易实现，却可能删除少见但关键的失败边界；按引用频率保留，又可能让经常出现的错误建议永久占据空间。使用频率反映流量，不直接等于因果价值或风险价值。

可以把评估问题写得更具体：删除某条经验后，哪些任务的行为可能改变？这种改变的收益、损失与维护成本如何比较？若没有足够数据估计单条贡献，可以按内容类别、风险和证据强度采用明确规则，并记录规则的适用范围。不要给每条经验编造一个看似精确、实际没有依据的价值分。

主动遗忘也有实验成本。为了判断一条经验是否有用而反复调用模型，可能比保留它更贵；为了降低提示长度做大量压缩，可能引入新的错误。维护策略本身需要计费与评估。第 13 章会把这些成本放进完整生命周期，本章先要求它们不被“后台整理”这个名称排除在账外。


\section{来源链：一条经验必须能回答它从哪来}
\label{sec:ch05-provenance}

筛选与抽象常从写入开始，合并与冲突处理也可能在后续维护中发生，
而一条经验被质疑通常发生在很久之后——
换了一个读取者、换了一版依赖、或者它在某道题上第一次变成有害的。
要让那时候的质疑有对象可查，经验条目必须带一条来源链。

\begin{table}[htbp]
\centering
\caption{来源链的四个字段：没有它会失去什么}
\label{tab:ch05-provenance}
\begin{tabular}{>{\raggedright\arraybackslash}p{3.0cm}>{\raggedright\arraybackslash}p{5.2cm}>{\raggedright\arraybackslash}p{4.4cm}}
\toprule
字段 & 缺了它，什么判断做不了 & 本书实验台在哪里记了它 \\
\midrule
写入者（模型名、温度、提示版本） & 缺少生成条件，难以回查错误是在提炼还是使用时引入 & 配置记录模型名与温度；提示版本还须核对源码，旧归档未绑定运行时源码哈希 \\
来源运行（运行号、来源任务、该次验收结果） & 验收结果无法回溯：不知道这条经验是从一次通过还是一次失败里提炼的 & 记忆文件的 \texttt{source\_task} / \texttt{source\_verdict}，结合所属运行目录与配置中的 \texttt{run\_id} \\
时间（时间戳与代码版本） & 过期判定没有基准：不知道这条经验是在环境哪一版上成立的 & \texttt{config.json} 的 \texttt{git\_head}；轨迹的 \texttt{started\_at} \\
是否被真正采用 & 无法区分``检索到了''与``用上了''，检索命中率会被当成效果 & \textbf{本书尚未记录；缺口与影响见本节正文} \\
\bottomrule
\end{tabular}
\end{table}

\begin{principlebox}
\textbf{来源链让撤销有明确对象，但不会自动判断对错。}
可以依据运行、候选版本与条目标识追踪一条经验怎样产生，再查其后续派生项。
时间戳只能提示核查新旧，不能单凭时间长短判定过期；写入模型相同也不保证当前读取者适用。
识别错误仍需检查事实、条件或效果，来源链提供的是回查与撤销的路径。

\end{principlebox}

\begin{sloppypar}
表~\ref{tab:ch05-provenance} 的第四行是本书自己的缺口，不是别人的：
本书的实验台记录了每条经验的来源任务与来源验收结果
（E03 的记忆文件里 \texttt{source\_task} 与 \texttt{source\_verdict} 两个字段），
但\textbf{没有记录这条经验被检索到之后是否被模型真正采用}。
因此本书能说``检索到了几条''，不能说``其中几条起了作用''，
不能对单条经验作确定的效果归因。整体对照收益仍可估计；使用行为与条目归因的缺口已登记进 \texttt{agent/BLOCKERS.md}。
\end{sloppypar}

\section{检索条件、适用范围与虚假相关}
\label{sec:ch05-retrieve}

虚假相关那篇工作把问题分成三类：显式混淆、未测量混淆、对撞偏倚\cite{trap-trajectory-2605}。
三类不是一个难度梯度，而是三种不同的因果结构，
因此``用一个方法压住三类''这件事本身有信息量。

它给出的处置办法（CAMEL）在两处干预\cite{trap-trajectory-2605}：
\begin{itemize}
  \item \textbf{写入时}：在其表示假设下减去同一步内共享的分量，并用不依赖下游结果的新颖性准则决定写入；
  \item \textbf{检索时}：检验一个候选的相关性在沿非因果方向扰动后是否稳定，
        把那些排名依赖于共享虚假结构的候选降权。
\end{itemize}

这两处校准不增加大语言模型的前向调用，也不需要微调\cite{trap-trajectory-2605}。
缺少节点嵌入时仍需编码，在线统计与重排序也有计算开销。原文在给定表示与干预假设下构造可计算的稳定性检查，
并在其基准上评估，不能由该统计量直接判定任意相关性是否具有因果支持。

\begin{sloppypar}
结合上述风险，本书采用一条写作与工程判据：
\textbf{有决策作用的经验应明确适用条件}。
KSI 的知识条目结构列出 \texttt{applies\_when} 与
\texttt{does\_not\_apply\_when}\cite{ksi-2607}。
字段提供一个可检查的边界，但不会自动证明条件正确；虚假相关仍需按具体机制检验。
\end{sloppypar}

\section{过期经验、淘汰与主动遗忘}
\label{sec:ch05-forget}

不同形态依赖不同条件：事实可能随世界变化，程序可能随接口变化，策略也可能因目标改变而不再适用。
章首案例涉及截断与检索，不是一个已经观察到的接口过期事故。

从记忆组织走向淘汰，还需要另一组明确操作：什么信号触发停用、删除哪些对象、怎样处理派生物，以及误删后能否恢复。
A-Mem 的笔记与演化机制提供组织经验的例子\cite{amem-2502}，不能仅凭这些机制推断淘汰策略已经得到验证。
本书E03冻结经验后进行评价，没有独立比较主动遗忘；若要评价淘汰策略，需要把规则和维护成本一并纳入新的比较。

一个不需要删除原语也能部分解决的替代办法，是 KSI 的 \texttt{does\_not\_apply\_when}：
不删，但把``什么时候不该用它''写进条目本身\cite{ksi-2607}。
这能保留已知边界，但是否降低误用取决于读取者是否检查条件。
条目数量、检索噪声与维护成本还受合并和索引策略影响，不能仅从字段形式推断。

\section{从同任务重用到跨任务迁移}
\label{sec:ch05-migrate}

\begin{sloppypar}
任务目录的划分与一次运行实际接触的任务，需要分别核对。
任务集 \texttt{v1} 的四个用途划分（train / select-valid / regression / sealed）
各 \calcnum{16} 道题，每个划分里 8 个族各 2 道；这只说明任务池的覆盖。
E03拆分B实际从前\measured{8}道训练题构建经验，覆盖\measured{4}族；
封存集的\measured{16}题则覆盖\measured{8}族。因而本轮整体读数混合了实际来源已见族与未见族，
既不是纯同族迁移估计，也不是纯跨族迁移估计。若要分别回答两类问题，须按来源暴露界定并分别分析，不能只依据train与sealed这两个目录名。
拆分A实际采集到的来源族与拆分B相同，候选训练池中还有其他族，却没有在本次来源回合配额内被采到。
\end{sloppypar}

\begin{principlebox}
\textbf{这条限制约束 P3 的可判定范围。}
E03 拆分 A 留出族名排序的最后两个族，所有训练候选都不含它们，测试取这两个族的全部任务。
这与标准划分不同，确实形成了相对于本次经验构建的族留出。
不过它只有两个留出族，也没有形式与容量的交叉操作；当前数据不足以回答“形式比容量更重要”。
区间跨零同样不能被解释为形式等价或命题反证。

\end{principlebox}

% ---------- 对照实验：E03 / E03c ----------
\section{两种拟合目标}
\label{sec:ch05-fit}

% figures/ch05-fit.tex —— 第 5 章核心图（全书引用入口）
% 用法：在 book/chapters/ch05.tex 中 % figures/ch05-fit.tex —— 第 5 章核心图（全书引用入口）
% 用法：在 book/chapters/ch05.tex 中 % figures/ch05-fit.tex —— 第 5 章核心图（全书引用入口）
% 用法：在 book/chapters/ch05.tex 中 \input{figures/ch05-fit}
% 图主体在 ch05-fit-body.tex，便于单独编译与复用。
\begin{figure}[htbp]
\centering
\input{figures/ch05-fit-body}
\caption{同一条来源记录可向求解器或问题域两种方向渲染，再交给同模型与换模型读取者。
图示为比较结构，不预设哪种形式更好。每个读取者都需要自己的无记忆基线；
还须核对渲染造成的内容、长度与检索变化，不能把同源等同于其他条件完全相同。}
\label{fig:ch05-fit}
\end{figure}

% 图主体在 ch05-fit-body.tex，便于单独编译与复用。
\begin{figure}[htbp]
\centering
% A comparison design, not an ordering of measured outcomes.
\begin{tikzpicture}[
  box/.style={draw,rounded corners=2pt,text width=26mm,minimum height=12mm,align=center,font=\small},
  note/.style={anchor=west,text width=30mm,align=left,font=\footnotesize},
  arrow/.style={-{Stealth[length=1.6mm]},thick}
]
\node[box,fill=gray!10] (s) at (0,0) {同一来源记录\\轨迹与验收};
\node[box,fill=orange!12] (a) at (4,1) {求解器方向\\候选行为补丁};
\node[box,fill=blue!10] (b) at (4,-1) {问题域方向\\带条件的主张};
\node[box,fill=green!10] (r1) at (8,1) {同模型读取};
\node[box,fill=green!10] (r2) at (8,-1) {换模型读取};
\draw[arrow] (s)--(a);
\draw[arrow] (s)--(b);
\draw[arrow] (a)--(r1);
\draw[arrow] (a)--(r2);
\draw[arrow] (b)--(r1);
\draw[arrow] (b)--(r2);
\node[note] at (9.7,1) {分别比较两种形式；\\减同模型无记忆基线};
\node[note] at (9.7,-1) {分别比较两种形式；\\减换模型无记忆基线};
\node[align=left,text width=140mm,font=\footnotesize] at (5.8,-2.2)
{同源固定材料来源，渲染仍可改变内容、长度与检索。图中箭头表示比较路径，不表示已证实的迁移收益。};
\end{tikzpicture}

\caption{同一条来源记录可向求解器或问题域两种方向渲染，再交给同模型与换模型读取者。
图示为比较结构，不预设哪种形式更好。每个读取者都需要自己的无记忆基线；
还须核对渲染造成的内容、长度与检索变化，不能把同源等同于其他条件完全相同。}
\label{fig:ch05-fit}
\end{figure}

% 图主体在 ch05-fit-body.tex，便于单独编译与复用。
\begin{figure}[htbp]
\centering
% A comparison design, not an ordering of measured outcomes.
\begin{tikzpicture}[
  box/.style={draw,rounded corners=2pt,text width=26mm,minimum height=12mm,align=center,font=\small},
  note/.style={anchor=west,text width=30mm,align=left,font=\footnotesize},
  arrow/.style={-{Stealth[length=1.6mm]},thick}
]
\node[box,fill=gray!10] (s) at (0,0) {同一来源记录\\轨迹与验收};
\node[box,fill=orange!12] (a) at (4,1) {求解器方向\\候选行为补丁};
\node[box,fill=blue!10] (b) at (4,-1) {问题域方向\\带条件的主张};
\node[box,fill=green!10] (r1) at (8,1) {同模型读取};
\node[box,fill=green!10] (r2) at (8,-1) {换模型读取};
\draw[arrow] (s)--(a);
\draw[arrow] (s)--(b);
\draw[arrow] (a)--(r1);
\draw[arrow] (a)--(r2);
\draw[arrow] (b)--(r1);
\draw[arrow] (b)--(r2);
\node[note] at (9.7,1) {分别比较两种形式；\\减同模型无记忆基线};
\node[note] at (9.7,-1) {分别比较两种形式；\\减换模型无记忆基线};
\node[align=left,text width=140mm,font=\footnotesize] at (5.8,-2.2)
{同源固定材料来源，渲染仍可改变内容、长度与检索。图中箭头表示比较路径，不表示已证实的迁移收益。};
\end{tikzpicture}

\caption{同一条来源记录可向求解器或问题域两种方向渲染，再交给同模型与换模型读取者。
图示为比较结构，不预设哪种形式更好。每个读取者都需要自己的无记忆基线；
还须核对渲染造成的内容、长度与检索变化，不能把同源等同于其他条件完全相同。}
\label{fig:ch05-fit}
\end{figure}


一条经验可以往两个方向写：
\begin{itemize}
  \item \textbf{拟合到求解器}：写成``针对当前这个模型/当前这个 harness 的补丁''。
        它可以依赖这个模型的具体习惯、措辞与默认行为。
  \item \textbf{拟合到问题域}：写成``与谁读它无关的事实或不变量''，
        并带上适用与不适用条件。
\end{itemize}

这是一个判断，不是既有结论。
\texttt{spec/glossary.md}将它登记为本书采用的区分；E03c提供当前限定条件下的比较，
尚未裁定更广的迁移主张，不能把这一区分写成既有定律或已经核实的原创性结论。
两侧各有一个单边数据点：
harness 进化在其测试划分上的平均增益 \cited{+0.6} 个百分点（同一报告里 GPT-5.4 为 \cited{+0.0}）\cite{harness-eval-2607}；
冻结知识束在作者选定的未见任务上提供跨模型族收益证据\cite{ksi-2607}。
两个数据点来自不同任务、不同协议，\emph{不能直接相减}，
这不构成``问题域拟合更好''的证据。

\begin{sloppypar}
A-Mem 在局限一节指出，组织质量可能受底层语言模型能力影响，
不同模型可能生成不同描述与链接\cite{amem-2502}。
这里首先说的是\emph{构建者依赖}：换构建模型，产物可能变化。
它还没有证明\emph{读取者依赖}：固定同一产物，换读取模型，收益会怎样变化。
即使一项实现恰好由同一模型承担两种角色，也不能把两种问题合并。
前者说明有必要检查生成过程，后者才直接涉及冻结经验的可移植性；
二者都不能仅从“内容由模型生成”推出经验必然拟合了某个求解器的弱点。
\end{sloppypar}

\begin{principlebox}
本轮 E03 没有独立操纵记忆容量，因此不能检验“形式比容量更重要”的完整 P3。
拆分 A 只有两个留出族，题级区间仅描述当前固定族中的比较，不足以支持任意新族的迁移主张。
E03的十臂（E03四臂 + E03c六臂）均有执行归档，E03c最后一臂于2026-09-28续跑补齐；拆分A基线因记录冲突暂停采用，见下文。
矩阵完整不消除跨日期模型别名、经验选择与统计精度的限制，拟合目标与读取者变化的交互仍需新的对照。
\end{principlebox}

% ==E03_RESULTS_START==
\subsection{本书的实测：E03 与 E03c}
\label{sec:ch05-e03}

\begin{sloppypar}
运行 \texttt{E03-20260925T152453Z-254845}：同模型 \texttt{deepseek-flash}、换模型 \texttt{deepseek-v4-pro}，训练集配额 $E=$\measured{8}，每题重启回合 $R=$\measured{5}，留出族 \texttt{sortkey} 与 \texttt{strindex}，协议所用 MDE 设计值 \calcnum{0.03}，重采样单位为题。两套拆分见 §\ref{sec:ch05-migrate}。
\end{sloppypar}

\subsubsection{拆分 A（族留出）：P3 的相关比较}

\begin{table}[htbp]
\centering
\caption{拆分 A（族留出）：可追溯的记忆臂与暂停使用的基线}
\label{tab:ch05-e03a}
\small\setlength{\tabcolsep}{4pt}
\begin{tabular}{>{\raggedright\arraybackslash}p{3.3cm}cc>{\raggedright\arraybackslash}p{4.7cm}}
\toprule
臂 & 回合通过率 & 通过 / 总运行 & 本次采用范围 \\
\midrule
\texttt{base\_same} & 暂停采用 & 记录冲突 & 汇总与保存轨迹不一致 \\
\texttt{form\_raw} & \measured{0.8875} & \measured{71}/\measured{80} & 描述该臂；相对基线收益暂停 \\
\texttt{form\_summary} & \measured{0.7750} & \measured{62}/\measured{80} & 描述该臂；相对基线收益暂停 \\
\texttt{form\_structured} & \measured{0.7625} & \measured{61}/\measured{80} & 描述该臂；相对基线收益暂停 \\
\bottomrule
\end{tabular}

\vspace{2mm}
\begin{footnotesize}
运行\texttt{E03-20260925T152453Z-254845}，留出族共\measured{16}题、每题\measured{5}次。三个记忆臂的逐次判定及用量与保存轨迹相符。基线的历史汇总和检查点记\measured{72/80}，保存轨迹的最终判定合计为\measured{70/80}；差异集中在\texttt{sortkey-03}的r2、r3。二者不是两份可互换的有效结果。
\end{footnotesize}
\end{table}

\paragraph{审计案例：先确定记录集，再计算效应。}
\label{sec:ch05-record-conflict}
旧 E03 运行器将轨迹判定写入逐题向量，再将向量保存在检查点；续跑按已完成臂恢复该向量。留存的旧检查点没有把每个判定绑定到轨迹哈希，也没有保存可解释这两次差异的替换链。首次Git归档中已经存在同一冲突，不能用后来的复制或舍入解释。修订版不在两种计数之间挑选一个作为真值：暂停该基线、三个相对基线的点差与区间，以及依赖它的成本效率解释。“三种记忆形式均低于基线”的概括随之撤回。

用保存轨迹直接换算得到\calcnum{0.8750}，但这只是冲突一侧的算术；它没有证明该侧就是原分析采用的有效记录。基线用量也不一致：保存轨迹合计比历史臂账少\calcnum{8}次调用、\calcnum{3839}个提示token及\calcnum{351}个生成token。附录~\ref{app:cost}因此将原账标为历史汇总，不据此作已校正成本结论。完整逐次身份和差异保存在交付包的analysis目录：

\begin{itemize}
  \item \nolinkurl{analysis/E03-record-index.csv}：逐次轨迹与哈希。
  \item \nolinkurl{analysis/E03-provenance-audit.json}：判定、用量与汇总链差异。
\end{itemize}

这一案例确定了数据使用边界：来源冲突影响的是该基线及其派生量。三个记忆臂之间的直接比较不使用该基线，可以继续核对；恢复基线后，形式与容量的交叉设计仍须另行建立。后文涉及此基线时统一返回本节。

历史协议为 P3 选取的局部比较是结构化减摘要这一对配对差（表~\ref{tab:ch05-e03a}）：

\begin{itemize}
  \item 点估计 \measured{-0.0125}，95\% 区间 \measured{[-0.1500, +0.1125]}，配对题数 \measured{16}。
  \item 区间包含 0：本轮未能确定结构化与摘要差异的方向，\emph{这不构成等价或证伪}；P3 还缺容量对照。
  \item 参照：原始轨迹减摘要，点估计 \measured{+0.1125}，区间 \measured{[-0.0125, +0.2375]}。
\end{itemize}

\subsubsection{拆分 B（标准划分）：P3b 的相关比较}

\begin{table}[htbp]
\centering
\caption{拆分 B（标准划分）：两种拟合目标 × 两个读取者}
\label{tab:ch05-e03c}
\small\setlength{\tabcolsep}{4pt}
\begin{tabular}{>{\raggedright\arraybackslash}p{3.2cm}ccc}
\toprule
臂 & 回合通过率（题均） & 通过 / 总运行 & 相对同读取者基线的收益（95\% 区间） \\
\midrule
\texttt{base\_same} & \measured{0.6875} & \measured{55}/\measured{80} & （基线） \\
\texttt{base\_cross} & \measured{0.2625} & \measured{21}/\measured{80} & （基线） \\
\texttt{fit\_solver\_same} & \measured{0.7000} & \measured{56}/\measured{80} & \measured{+0.0125}（\measured{[-0.1125, +0.1250]}） \\
\texttt{fit\_solver\_cross} & \measured{0.3625} & \measured{29}/\measured{80} & \measured{+0.1000}（\measured{[-0.0125, +0.2125]}） \\
\texttt{fit\_domain\_same} & \measured{0.7000} & \measured{56}/\measured{80} & \measured{+0.0125}（\measured{[-0.1500, +0.1625]}） \\
\texttt{fit\_domain\_cross} & \measured{0.4250} & \measured{34}/\measured{80} & \measured{+0.1625}（\measured{[+0.0000, +0.3250]}） \\
\bottomrule
\end{tabular}

\vspace{2mm}
\begin{footnotesize}
数字身份：[实测] 本书运行 \texttt{E03-20260925T152453Z-254845}；判决集为 sealed 划分 \measured{16} 道。同模型臂减 \texttt{base\_same}，换模型臂减 \texttt{base\_cross}；直接相减不同读取者的总分不能单独表示记忆收益，本表使用各自基线。
\texttt{fit\_domain\_cross} 一臂于 2026-09-28 用 \texttt{--resume-latest} 续跑补齐（原阻塞于换模型后端余额，HTTP 402）；收益区间的下界显示为 \measured{+0.0000} 是浮点舍入的结果，
实际下界与 0 重合，\textbf{不得读作``下界为正''}。
\end{footnotesize}
\end{table}

历史协议为 P3b 选取的局部比较是换模型后问题域减求解器这一对配对差：

\begin{itemize}
  \item 换模型后：点估计 \measured{+0.0625}，95\% 区间 \measured{[-0.0625, +0.1750]}，配对题数 \measured{16}。
  \item 区间包含 0：\emph{这不构成等价或证伪}——两种拟合目标换模型后的差异方向在本轮样本下仍不确定，拟合目标差异仍需在更广读取者范围中检验。
  \item 同模型参照：点估计 \measured{0.0000}，区间 \measured{[-0.1375, +0.1500]}。
\end{itemize}

逐题收益点估计反号的题数（同读取者下高于各自基线、换读取者后低于各自基线）：

\begin{itemize}
  \item \texttt{fit\_solver}：\measured{1} / \measured{16} 道题反号。
  \item \texttt{fit\_domain}：\measured{2} / \measured{16} 道题反号。
\end{itemize}

\subsubsection{历史精度估算与设计缺口}

\begin{sloppypar}
配对差标准差：P3 \measured{0.2778}、P3b \measured{0.2500}。按 $n\ge(1.96\,s/h)^{2}$ 作历史区间半宽规划，令 $h=$ \calcnum{0.03} 需要 P3 \calcnum{330} 道配对题、P3b \calcnum{267} 道配对题；当前可用为拆分 A \measured{16} 道、拆分 B \measured{16} 道。
所需规模约为当前规模的\calcnum{20.6}倍（P3）与\calcnum{16.7}倍（P3b）。
\textbf{该式是精度近似而非功效设计}，确认性阶段须按具体目标完成样本规划。
\end{sloppypar}

% ==E03_RESULTS_END==

\section{可移植性不是一个分数，而是一组变化条件}
\label{sec:ch05-portability}

经验的可移植性至少涉及任务、读取者、执行工具与环境。保持模型不变、换一道同族题，检验的是题目变化；把经验交给另一模型，检验的是读取者变化；换工具协议或动作预算，检验的是执行约束变化。这些变化可以同时发生，但同时变化越多，越难把效果归到某一个因素。

\subsection{先问经验在为谁解决什么问题}

针对求解器的经验可以提醒某个执行流程容易忽略的步骤，针对问题域的经验可以描述与具体执行者无关的条件。两种方向都可能有价值：前者贴近当前系统的弱点，后者更容易表达为可复查的任务知识。但把文本写成第三人称，不会自动消除模型特定偏差；使用第二人称，也不表示其中没有普遍知识。

E03 拆分 B 的一条求解器经验对计数器任务说：“你容易把简单实现题扩成完整 unittest，连 reload/fresh 样板都带上；只改 main.py 实现 Counter，别写验收测试。”对应的来源轨迹却只记录了一次对源文件的正确实现写入，验收也通过。来源材料包含测试文件内容，但可见动作没有把整套测试写进提交。

因此，这条文字对模型习惯的诊断缺少该轨迹的直接支持。不能因来源材料中出现测试框架，就认定提交者复制了测试框架；也不能从一次任务可靠估计模型的稳定习惯。提示要求必须指出“你容易犯的错”，可能推动渲染器在没有足够失败证据时仍给出人格化的补丁。这是从实际归档文字发现的内容风险，不是对所有求解器经验的否定。

对应的问题域条目描述计数器状态、累加操作及适用条件，更贴近该任务的契约。不过其中关于并发、外部修改或非数值输入的边界，也应区分是已测反例还是基于规格提出的限制。写成“事实、适用、不适用”三行，只让这些主张容易找到，并没有替代核验。

若要改进渲染器，可以允许“没有足够证据识别特定习惯”，也可以将观察到的错误与一般提醒分别输出。这样得到的条目可能不如一句肯定诊断简洁，却避免把一个来源中的偶然现象固化为读取者身份的一部分。渲染过程本身也应接受内容检查，而不是只验字段是否齐全。

\subsection{换读取者时，用各自基线定义增益}

设读取者为 $r$，经验形式为 $m$，相应表现为 $S(m,r)$，无经验基线为 $S(0,r)$。一种可解释的比较是
\begin{equation}
G(m,r)=S(m,r)-S(0,r).
\label{eq:ch05-reader-gain}
\end{equation}
它把经验带来的变化与读取者自身的基础水平分开。若两个模型的原始成绩不同，不能直接把携带经验后的总分差称为经验可移植性的差异。

在同一个读取者下，比较问题域与求解器经验，可以直接用两者得分差，因为相同基线相消。但若问“换读取者以后，两种经验的相对优势是否改变”，需要比较两个读取者下的差异，而不只是看跨模型那一列。这涉及交互关系：
\begin{equation}
I=\bigl[G(\mathrm{domain},r_2)-G(\mathrm{solver},r_2)\bigr]
 -\bigl[G(\mathrm{domain},r_1)-G(\mathrm{solver},r_1)\bigr].
\label{eq:ch05-reader-interaction}
\end{equation}
式中的每个量都应来自对齐的任务、资源与数据条件，并共同估计不确定性；不能把两张表里分别显著与不显著的结果当作交互检验。

E03c 的完整矩阵为这种思考提供了起点，但跨日期续跑与模型别名的限制仍在。只看到某一读取者下的问题域点估计更高，尚不足以证明拟合目标决定所有迁移表现。本批两个模型别名、固定工具协议和小型任务集，限定了能讨论的范围。

\subsection{逐题反号用于定位，不自动建立条件规律}

历史报告把同读取者下高于基线、换读取者后低于基线的题记为反号。这个计数有助于找需要进一步检查的案例；但每题只由少量回合估计通过率，微小波动也可能跨过零线。把某题标成反号，不代表我们已经可靠证明经验在两种读取者下分别有益与有害。

下一步可以查看实际检索条目、提示中的规则冲突、工具格式和最终提交，提出解释，再用新的受控比较检验。若观察到格式失败集中在某读取者，可能首先需要区分接口适配与知识迁移；若两边都执行了相同动作但契约理解不同，再讨论经验表达。定位工作不能由一个计数替代。

还应避免只展示最戏剧化的反号题。完整分布包括改善、退化、持平与尚不确定的部分。挑选案例适合解释机制，整体结论应回到预定统计范围。案例可以告诉读者一种失效如何发生，却不能单独告诉读者它有多普遍。

\subsection{外部迁移证据要连同选择规则一起读}

KSI 将知识资产冻结后交给新的读取者，在未见任务上评估，不再进行新的论坛讨论或接收端蒸馏\cite{ksi-2607}。这比只展示同一运行越来越高的分数更接近持久物的可复用性问题。不过原文的任务选择与适配过程也必须一起保留。

其迁移实验在每个迁移基准中，从未参与自改进的候选题里，先寻找两种接收模型在一次预备无知识运行中都失败的题，再从交集中抽取 \cited{20} 题；测试时还使用面向当前任务的适配器，将冻结知识资产转成短备忘录\cite{ksi-2607}。因此，结果对应这组被筛选的困难题和该适配流程，不是任意新题上的无条件平均收益，也不是原始知识束未经处理的直接粘贴效果。
原文还指出，Polyglot迁移集中有\cited{7/20}道题与供体共享同一练习的不同语言版本\cite{ksi-2607}。任务标识留出与全新练习留出检验的范围不同，报告时需保留这一区别。

这并不使结果无效。它明确回答了“在这些两种读取者都曾失败的题上，携带冻结资产是否有帮助”这一问题。迁移资产、接收端适配器和题目筛选共同构成实验对象。阅读其他工作时，同样应问被冻结的到底是什么，接收端还做了哪些计算，评估题如何进入样本，以及是否有同条件的无资产基线。

本书的 E03c 与这一外部结果也不能直接相减。任务类型、读取者、知识构建过程、预算和选择条件都不同。外部证据可以反对过强的“记忆不能迁移”断言，也可以启发新设计，但不能替代本书矩阵中的未知参数。

\section{形式与容量的比较，怎样避免互相替代}
\label{sec:ch05-form-capacity}

P3关心形式与容量怎样共同影响迁移，而旧 E03 只提供固定条目数下几种完整表示方案的比较。原始截取、摘要与结构化长度不同，检索来源也经常不同。要继续推进，应先把容量和形式各自变成可操作的变量，而不是在现有表里寻找一个足够大的差值。

\subsection{给容量选择一个与问题一致的单位}

若关心存储维护成本，可以用字节、条目数量与索引开销；若关心推理时负担，更直接的是每次注入的 token 预算；若关心经验覆盖，可以用独立来源任务或条件类别的数量。它们回答不同问题。统一使用“容量”而不写单位，会让一个方案在条目数上更大、在提示长度上更小，却被随意解释为支持任意结论。

形式也应定义清楚。若一臂是自然语言摘要，另一臂是在相同文本前加字段名，那么研究的是格式提示；若另一臂增加反例和适用条件，研究的就是内容结构的联合变化。后者可能更符合工程需求，但不能把收益全部归给换行或字段名。

一种设计可以在同一来源集合上生成若干形式，再为每种形式设置多个固定的提示预算。若采用截断，应记录哪些部分被截掉；若采用摘要压缩，应把压缩过程与成本也纳入方案。不同形式不必共享完全相同的最优策略，但比较的目标必须从一开始说明，是固定规则下的公平比较，还是各自优化后的系统比较。

\subsection{比较形式效应与容量效应，需要固定参照范围}

“谁更重要”不是一个脱离范围的量。容量从很少增到适中，可能与从适中增到极大产生不同变化；形式从完整诊断改成关键词，也不同于从短段落改成三行字段。预先规定变化幅度与目标分布，才能比较这两个操作的实际影响。

可以分别估计在某个容量档位切换形式的差异，以及在某个形式下增加容量的差异，再检查两者是否交互。若结构化格式只有在较长预算下才保留完整条件，它的优势可能随容量改变。只报一个总体平均差，会隐藏这个关系；只报最有利的格点，又会夸大可推广范围。

容量还可能改变检索覆盖。多存条目不一定多读条目，多读条目也不一定多读不同来源；重复经验可能占据新增名额。应同时报告存储范围、候选集合、实际选中集合和最终展示长度。这样某个格点效果变化时，至少知道新增资源进入了哪一环。

\begin{table}[htbp]
\centering
\caption{形式与容量的交互：教学构造的得分，不是模型实测}
\label{tab:ch05-form-capacity-example}
\begin{tabular}{lcc}
\toprule
表示 & 低注入预算 & 高注入预算 \\
\midrule
摘要 & \calcnum{0.60} & \calcnum{0.70} \\
结构化 & \calcnum{0.65} & \calcnum{0.85} \\
\bottomrule
\end{tabular}
\end{table}

表~\ref{tab:ch05-form-capacity-example} 给出一个人为设定的例子，得分越高越好。在低预算下，结构化减摘要为\calcnum{0.05}；在高预算下，这个差值为\calcnum{0.15}。两者之差为
\begin{equation}
I_{\rm form,capacity}
=(0.85-0.70)-(0.65-0.60)=\calcnum{0.10}.
\label{eq:ch05-form-capacity-interaction}
\end{equation}
它描述这两个预算档位之间的形式差异怎样变化，不是某种格式在任意容量下的固定收益。反过来，扩容在摘要与结构化下分别增加\calcnum{0.10}与\calcnum{0.20}，也没有一个脱离表示的统一容量效应。若要比较“谁更重要”，还必须先规定变化范围与汇总权重。真实研究应对同一交互对比共同估计不确定性；一个格点显著、另一个不显著，不能代替交互检验。

\subsection{将构建成本与重复使用次数一起报告}

经验生成通常有前期成本，读取与维护也会持续消耗资源。若只报告测试阶段的调用，复杂渲染器看起来可能与直接截取一样便宜；若把构建成本全部压到一次使用，又可能低估长期复用价值。需要明确成本分摊的使用次数与任务分布，并允许某些资产在实际流量下始终不能回本。

E03 的账本提供了一次具体教训。最终汇总中的 \texttt{side\_ledger} 为零，但来源轨迹显示早期确实完成了 \measured{16} 个回合、\measured{79} 次调用，包含 \measured{29443} 个输入 token 与 \measured{4243} 个输出 token。续跑恢复了已完成测试臂，却没有从中断记录恢复来源与渲染总账。

按保留渲染实现与已保存条目，可以重建有 \calcnum{48} 次成功的模型渲染调用。这个总数还包括拆分B已渲染、却没有设相应评价臂的摘要与结构化条目，分摊成本时需要保留这部分构建活动的范围。这些调用的原始 token 用量没有归档，不能从最终经验字数反推。因而报告应区分恢复出的实测来源成本、由保存条目重建的调用数，以及仍然缺失的渲染用量。

这类缺口不会把已经保存的作答成绩变成不存在，却会限制资源效率结论。如果不知道构建成本，就不能可靠声称某种形式更省；如果只知道调用次数，也不能把所有调用当作等价。成本缺失与成绩缺失是不同问题，应分别列出，而不是用一个“运行完成”状态掩盖它们。

\subsection{小规模预检应验证什么}

在扩大规模以前，可以离线检查条目是否包含预定字段、是否超限、是否混入不允许的任务材料、检索是否使用了预定源条目，以及提示是否真正携带这些内容。用固定的虚拟后端还可以检查成本累加与中断恢复，不必先付费生成大量结果。

内容检查则需要抽取有代表性的成功与失败来源，核对摘要是否保留关键条件、是否把测试样板误认作提交代码、是否为成功任务编造模型缺陷。离线规则只能发现一部分问题；人工或模型审查也只能为已审范围提供证据，应保留具体审查对象与结论。

工程预检通过以后，才适合冻结新版本并运行预定比较。预检期间发现问题而改动形式、预算或检索策略，是正常设计过程；一旦结果用于主要结论，就不能再把同一批数据当未经选择的确认样本。把开发、探索与确认分开，能让负结果指出真实边界，也能让正结果获得清楚的解释。


\section{把一条有问题的经验修订成可审查版本}
\label{sec:ch05-revision-card}

回到计数器那条求解器补丁，可以把一次完整维护过程写出来。下面的修订是本书根据现有轨迹提出的教学版本，没有写回旧实验的记忆文件，也没有让模型据此重新作答。原经验与原成绩保持原样，教学修订用于说明怎样处理不受来源支持的主张。

第一步是把原条目拆成两个主张。一个主张是“该模型容易把测试样板带进实现”，另一个是“本任务应修改源文件，不应写验收测试”。后者与当前任务的保护规则相容；前者涉及模型的一贯行为，仅凭这一条成功轨迹没有足够支持。不能因为后一半提醒合理，就连带认可前一半诊断。

第二步是回查动作。来源轨迹记录了写入 \texttt{main.py} 的计数器实现，内容只有初始化状态与累加返回逻辑。验收测试文本属于输入材料的一部分，出现在渲染器阅读的代码集合中。于是更贴近证据的解释是：渲染器可能混淆了“提供给它阅读的文件”和“提交者生成的文件”。这仍是一种根据输入与输出提出的机制假设，不能断言已经读取了渲染模型的内部推理。

第三步是保留可支持的部分，并改写证据状态。教学修订可写为：“本来源只修改了计数器实现并通过已知测试；未观察到生成测试样板的行为。后续同类任务应遵守测试不可写的规则；尚无充分证据为该模型归纳稳定的样板复制倾向。”它不再像一句锋利的补丁，却清楚区分观察、任务约束和未证主张。

第四步是决定放在哪里。这条修订更适合进入来源审计与渲染器回归用例，不一定值得作为每道编程题的经验反复注入。如果提示预算有限，源文件修改规则本来已在系统指令中出现，重复提醒可能没有新增价值。维护的结果可以是撤销原条目的自动使用、保留审计记录，并把新案例用于检查下一版经验生成器。

第五步是提出后续检查，而不是宣布问题已经彻底解决。可以用固定材料比较两种输入构造：一种混合源文件和测试文件，另一种明确标记初始文件、实际改动与受保护文件。检查输出是否仍会把测试内容归到提交者行为上。若误归因减少，说明这项输入设计在所测范围内有帮助；若没有减少，则需要继续检查提示是否强迫生成缺陷诊断。

这个案例说明，经验维护有时不会产出更长、更强的经验，而会产出一个更窄的结论和一个更好的实验。成熟的经验系统应当允许这样的结果。否则每次审计都被要求“提炼一条普遍规律”，错误主张就可能在修订过程中只是换一种更谨慎的措辞继续存活。

对发布者而言，还要留下两条连接：原条目指向为何停用的证据，新渲染器的检查用例指向原故障。这样之后再遇到相似输出，可以判断旧问题是否复发，也能防止有人只看到原经验语句流畅就重新启用。经验的价值不仅在于告诉执行者做什么，也在于让系统知道哪些说法曾被检查、被限制或被撤销。


% ---------- 5. 条件变化 ----------
\section{条件变化：换题、换族、换模型}
\label{sec:ch05-conditions}

\begin{itemize}
  \item \textbf{换题（同族新题）。}测试题所属的族须在本次经验来源中出现过。
        标准任务池覆盖同一组族，但实际训练配额未必覆盖全部族；
        E03拆分B的整体读数包含来源已见与未见族，须分开分析才可定位这一维。
  \item \textbf{换族。}测试族未进入本次经验来源。拆分A是专门设置的族留出比较（§\ref{sec:ch05-migrate}），
        与P3相关；其他划分也须按实际来源核对，不能只凭名称排除或确认换族。
  \item \textbf{换模型。}这一维是 P3b 的相关比较。
        关键设计是必须有一个\emph{换模型后的无记忆基线}，
        否则``换模型后收益下降''无法区分是经验不可移植，还是另一个模型本来就弱。
\end{itemize}

% ---------- 6. 选择原则 ----------
\section{选择原则}
\label{sec:ch05-choice}

\begin{enumerate}
  \item \textbf{决策性经验应写出已知边界。}
        字段为空应视为边界未充分说明；字段齐全也需要证据，不能把条件文字当作正确性证明。
  \item \textbf{报迁移数字时必须先说清换的是什么。}
        换题、换族、换模型是三个不同的量，混着报会让 P3 与 P3b 都无法判。
  \item \textbf{换模型的主张必须带换模型的静态基线。}
        这是 E03c 里 \texttt{base\_cross} 那一臂存在的全部理由。
  \item \textbf{两种不同的拟合目标要分别渲染，且必须同源。}
        同源固定的是材料来源；渲染仍可能改变内容、长度与检索，必须继续核对这些路径。
  \item \textbf{写不出删除原语时，至少把误用记下来。}
        \texttt{does\_not\_apply\_when} 不是删除的替代品，
        还需要来源、版本与实际采用记录来支持审计。
\end{enumerate}

% ---------- 7. 反例 ----------
\begin{pitfallbox}
\textbf{反例一：把``存下来了''当成``学会了''。}
第 3、4 章的部分对照中，带记忆的回合通过率低于相应无记忆臂。
存下来只是写入成功，不是能力变化；
两者之间隔着一整套检索与适用条件的检验。

\textbf{反例二：把检索到的相关性当成因果。}
虚假相关的两个性质（持久性、反馈）会把一个挥发性错误变成自我强化的偏置
\cite{trap-trajectory-2605}。
该论文在所测校准基线中报告了指标之间的残余权衡，不能外推为所有经验系统的必然规律。

\textbf{反例三：把同族新题上的提升当成跨族迁移。}
任务集 v1 的每个划分里 8 个族各 2 道（附录 B §\ref{sec:appb-tasks}），
标准划分名本身不保证跨族留出，还须核对实际来源与暴露；本章的族留出比较来自另设的拆分 A。

\textbf{反例四：把``换模型还能跑''当成``经验可移植''。}
判别方法是补一个换模型的无记忆基线。
A-Mem 自陈组织质量可能受底层模型能力影响\cite{amem-2502}，
这支持检查构建者差异，不能代替固定记忆后的读取者迁移对照。
\end{pitfallbox}

\section*{习题}
\begin{enumerate}[label=5-\arabic*.,ref=5-\arabic*]
  \item\label{ex:ch05-1} 取你正在用的记忆系统，按表~\ref{tab:ch05-forms} 给它的每条经验归类。
        哪一类最多？那一类的典型失效在你的系统里出现过吗？
  \item\label{ex:ch05-2} 式~\eqref{eq:ch05-link} 用余弦相似度取近邻。
        请设计一个反例，使得两条语义相近但因果无关的经验被排到前 $k$。
  \item\label{ex:ch05-3} 本章说``不删，但把不适用条件写进条目''是删除原语的部分替代。
        请列出它无法替代的三件事。
  \item\label{ex:ch05-4} E03c 的换模型臂要求同源渲染。
        若两臂各自由不同的轨迹生成，来源内容、质量与任务覆盖的可能变化会怎样影响比较？
\end{enumerate}

\section*{常见误解}
\begin{itemize}
  \item \textbf{``记忆越多越聪明''。}
        章首展示的是关键内容被截掉、渲染改变检索；它没有独立操纵容量。
        外部工作另展示记忆可能保存并反馈强化虚假线索\cite{trap-trajectory-2605}。
        两者提示不同风险，不能合成“加容量必然有害”的结论。
  \item \textbf{``压缩一下就好了''。}
        Mem-$\alpha$ 把压缩率写成一项奖励\cite{memalpha-2509}，
        其他条件相同时，缩短记忆可以提高该奖励分量；总目标还包括问答正确率、工具格式与内容合法性。
        第7章另讨论选择器偏爱短提示的问题，两种机制应分别检验是否丢失任务条件，不能仅凭文本变短判断压缩有效。
  \item \textbf{``检索到就说明相关''。}
        相关不等于有因果路径。三类虚假结构（显式混淆、未测量混淆、对撞偏倚）
        都会产生``检索得到''而没有因果支持的情况\cite{trap-trajectory-2605}。
  \item \textbf{``跨模型能用就是可移植''。}
        可能是另一个模型本来就强，或任务本身没区分度。
        先补换模型的静态基线。
\end{itemize}

\section*{进一步阅读}
\begin{itemize}
  \item \cite{trap-trajectory-2605}：持久性与反馈两条性质、
        三类虚假结构、以及 CAMEL 的两处干预。
        其中“每个基线”的残余权衡是对所测校准基线的讨论，不能扩大到所有记忆方法。
  \item \cite{amem-2502}：式~\eqref{eq:ch05-link} 与七字段笔记的出处，
        §4.3 是 \cited{1{,}200} token 与 \cited{0.0003} 美元的位置。
        按§4.4表3中实际移除的模块解释消融，避免把模块贡献扩大为纯拓扑贡献。
  \item \cite{ksi-2607}：\texttt{applies\_when} / \texttt{does\_not\_apply\_when} 的出处，
        以及在作者筛选的困难题与接收端适配流程下，冻结知识束的跨模型迁移结果。
  \item \cite{harness-eval-2607}：harness 进化迁移 \cited{+0.6}（GPT-5.4 为 \cited{+0.0}）
        的位置——拟合到求解器那一侧的单边数据点。
  \item \cite{raw-history-memory-2610}：固定已构建内容，改变检索器与预算，并以来源标注辅助区分构建和读取瓶颈；这是对话记忆的预印本研究。
  \item 本书附录 B §\ref{sec:appb-tasks}：任务集划分与``同族''限制的出处。
\end{itemize}
